\documentclass[10pt]{article}
\usepackage{braket}
\usepackage{algorithm}
\usepackage{algpseudocode}
\usepackage{soul}
\usepackage{float}
\usepackage{tabularx}
\usepackage{mathtools}
\usepackage{amsmath, amssymb, amsthm}
\usepackage{tikz}
\usetikzlibrary{arrows}
\usetikzlibrary{decorations.text}
\usepackage{enumitem}
\usepackage{geometry}
\usepackage{fancyhdr}
\usepackage{titlesec}
\usepackage{xltabular}
\usepackage{color}
\usepackage{hyperref}
\geometry{margin=1in}

\pagestyle{fancy}
\fancyhf{}
\cfoot{\thepage}
\title{SFWRENG 3DX4 - Pre-lab 1}
\author{Matthew Farah, \href{mailto:farahm11@mcmaster.ca}{$\langle$farahm11@mcmaster.ca$\rangle$}, 400468251}
\date{\today}

\hypersetup{
	colorlinks=true,
	linkcolor=blue,
	filecolor=magenta,
	urlcolor=blue,
	citecolor=blue,
	anchorcolor=blue
}

\fancyhead{}
\fancyhead[L]{Matthew Farah, \href{mailto:farahm11@mcmaster.ca}{$\langle$farahm11@mcmaster.ca$\rangle$}, 400468251}
\fancyhead[R]{SFWRENG 3DX4 - Pre-lab 1}

\AtBeginEnvironment{align}{\setcounter{equation}{0}}

\newcommand{\comment}[1]{\parbox[t]{0.45\textwidth}{\raggedright #1}}

\setlength{\parindent}{0cm}

\setcounter{tocdepth}{3}
\hypersetup{bookmarksopen=true, bookmarksopenlevel=2, bookmarksdepth=3}

\newcounter{question}
\renewcommand{\thequestion}{\arabic{question}}

\newenvironment{question}{
	\refstepcounter{question}
	\phantomsection
	\addcontentsline{toc}{section}{\protect{\large{Question \thequestion}}}
	\vspace{1em}
	\par
	\large\textbf{Question \thequestion}
	\vspace{.2cm}
	\par
	\setcounter{subquestion}{0}
	\setcounter{step}{0}
}{\vspace{.5em}}

\newenvironment{solution}{
	\vspace{1em}\par\large\textbf{{Solution:}}\par\vspace{1em}
}{
	\vspace{1em}
}

\newcounter{subquestion}
\renewcommand{\thesubquestion}{\alph{subquestion}}

\newenvironment{subquestion}{
	\refstepcounter{subquestion}
	\phantomsection
	\vspace{1em}\par\large\textbf{(\thesubquestion)}\hspace{-.5em}
	\setcounter{step}{0}
	\addcontentsline{toc}{subsection}{\protect{\large{Part \thequestion.\thesubquestion}}}
}{
	\par
	\vspace{1em}
}

\makeatother

\makeatletter

\newcounter{step}
\renewcommand{\thestep}{\Roman{step}}

\newenvironment{step}{
	\refstepcounter{step}
	\vspace{1.5em}\par\noindent\large\textbf{(\thestep)}
}{
	\par
	\vspace{1.5em}
}

\makeatother

\newcolumntype{Y}{>{\centering\arraybackslash}X}
\renewcommand{\arraystretch}{1.8}

\fontsize{10pt}{14pt}\selectfont

\begin{document}

\maketitle

\tableofcontents
\newpage

\begin{question}
	The general open-loop transfer function which models the angular velocity $\omega(t)$ of a motor is:

	\begin{gather*}
		G_{\omega} = \frac{\Omega(s)}{R(s)} = \frac{A}{\tau{s} + 1} \\
	\end{gather*}

	Where $A$ and $\tau$ are positive, real-valued constants and $\Omega(s) = \mathcal{L}\set{\omega(t)}$.
\end{question}

\begin{subquestion}
	What is the transfer function for the angular position of the motor $\theta(t)$?
\end{subquestion}

\begin{solution}
	We know that the angular velocity of a rotating body is given by:

	\begin{gather*}
		\omega(t) = \frac{d\theta(t)}{dt} \\
	\end{gather*}

	And since:

	\begin{gather*}
		L\set{\omega(t)} = \Omega(s) \implies L \left\{ \frac{d\theta(t)}{dt} \right\} = \Omega(s) \implies L\set{\theta(t)} = \frac{1}{s}\Omega(s) \\
	\end{gather*}

	We can express the transfer function in terms of $\theta(t)$ ($G_{\theta}$) as:

	\begin{align*}
		G_{\theta} &= \frac{L\set{\theta(t)}}{R(s)} \\
		&= \frac{1}{s} \cdot \frac{\Omega(s)}{R(s)} \\
		&= \frac{1}{s} \cdot \frac{A}{\tau{s} + 1} \\
		&= \frac{A}{\tau{s}^{2} + s}
	\end{align*}
\end{solution}

\begin{subquestion}
	What, if any, is the steady-state value of $\omega$ in open loop response to a step input:

	\begin{gather*}
		r(t) = U_{o}u(t) = \begin{cases} U_{o}, \quad t \ge 0\\ 0, \quad \; \; t < 0 \end{cases}
	\end{gather*}

	Where $r(t) = \mathcal{L}^{-1}\set{R(s)}$ and $U_{o} \in \mathbb{R}$ is a constant voltage.
\end{subquestion}

\begin{solution}
	We know that:

	\begin{align*}
		\frac{\Omega(s)}{R(s)} &= \frac{A}{\tau{s} + 1} \implies \\
		\Omega(s) &= \frac{A}{\tau{s} + 1}R(s) \\
	\end{align*}

	Thus, since $\omega(t) = \mathcal{L}^{-1}\set{\Omega(s)}$, to find the
	steady state response of $\omega(t)$ (its behaviour as $t \rightarrow
	\infty$), we must find $L^{-1}\set{\frac{A}{\tau{s} + 1}}R(s)$. Let's
	start by finding $R(s)$.

	\begin{align*}
		R(s) &= \mathcal{L}\set{r(t)} \\
		&= \int_{0}^{\infty}r(t)e^{-st}dt \\
		&= \int_{0}^{\infty}U_{o}e^{-st}dt \\
		&= U_{o}\int_{0}^{\infty}e^{-st}dt \\
		&= -\frac{1}{s}U_{o}e^{-st} \big|_{t = 0}^{t = \infty} \\
	\end{align*}

	And since $e^{-st} \rightarrow 0$ as $t \rightarrow \infty$:

	\begin{align*}
		R(s) &= \frac{1}{s}U_{o}e^{-s(0)}
		&= \frac{U_{o}}{s} \\
	\end{align*}

	Therefore:

	\begin{align*}
		\Omega(s) &= \frac{A}{\tau{s} + 1}R(s) \\
		&= \frac{A}{\tau{s} + 1}\left( \frac{U_{o}}{s} \right) \\
		&= \frac{AU_{o}}{s(\tau{s} + 1)}
	\end{align*}

	The next step to compute $\omega(t)$, requires us to perform partial
	fraction decomposition on $\frac{AU_{o}}{s(\tau{s} + 1)}$ so we can
	find it from a lookup table. \\

	Letting $B$ and $C$ be variables we will solve for, we can express $\frac{AU_{o}}{s(\tau{s} + 1)}$ as:

	\begin{align*}
		\frac{AU_{o}}{s(\tau{s} + 1)} &= \frac{B}{s} + \frac{C}{\tau{s} + 1} \implies \\
		AU_{o} &= B(\tau{s} + 1) Cs \\
		&= B\tau{s} + B + Cs \\
	\end{align*}

	By equating coefficients, we see that:
	
	\begin{gather*}
		B = AU_{o} \\
		B\tau{s} + C{s} = 0 \implies C = -B\tau \implies C = -AU_{o}\tau \\
	\end{gather*}

	And therefore:

	\begin{gather*}
		\frac{AU_{o}} = \frac{AU_{o}}{s} - \frac{AU_{o}\tau}{\tau{s} + 1} \\
	\end{gather*}

	Now, we can finally compute $w(t)$!

	\begin{align*}
		\omega(t) &= \mathcal{L}^{-1}\set{\Omega(s)} \\
		&= \mathcal{L}^{-1} \left\{ \frac{AU_{o}}{s} + \frac{-AU_{o}\tau}{\tau{s} + 1} \right\} \\
		&= \mathcal{L}^{-1} \left\{ \frac{AU_{o}}{s} \right\} + \mathcal{L}^{-1} \left\{ \frac{-AU_{o}\tau}{\tau{s} + 1} \right\} \\
		&= \mathcal{L}^{-1} \left\{ \frac{AU_{o}}{s} \right\} + \mathcal{L}^{-1} \left\{ \frac{-AU_{o}\tau}{\tau(s + 1/\tau)} \right\} \\
		&= \mathcal{L}^{-1} \left\{ \frac{AU_{o}}{s} \right\} + \mathcal{L}^{-1} \left\{ \frac{-AU_{o}}{s + 1/\tau} \right\} \\
		&= AU_{o}\mathcal{L}^{-1} \left\{ \frac{1}{s} \right\} - AU_{o}\mathcal{L}^{-1} \left\{ \frac{1}{s + 1/\tau} \right\} \\
	\end{align*}

	And by use of a lookup table, we can therefore conclude that:

	\begin{gather*}
		\omega(t) = AU_{o}u(t) - AU_{o}e^{-\frac{1}{\tau}t} \\
	\end{gather*}

	As stated previously, the steady state of $\omega(t)$ is its behaviour as $t \rightarrow \infty$, and we can see that:

	\begin{align*}
		\lim_{t \rightarrow \infty} \omega(t) &= \lim_{t \rightarrow \infty}AU_{o}u(t) - AU_{o}e^{-\frac{1}{\tau}t} \\
	\end{align*}

	Since $\tau > 0$, we know that $e^{-\frac{1}{\tau}t} \rightarrow 0$ as
	$t \rightarrow \infty$, and because $\forall t \ge 0, \; u(t) = 1$ trivially
	means $u(t) \rightarrow 1$ as $t \rightarrow \infty$ we can finally conclude that:

	\begin{align*}
		\lim_{t\rightarrow\infty} \omega(t) &= AU_{o}(1) - AU_{o}(0) \\
		&= AU_{o} \\
	\end{align*}

	Thus the steady state of $\omega(t)$ does indeed exist and is given by $AU_{o}$.
\end{solution}

\end{document}
